\input{../preamble.tex}

\lecturenumber{6}
\title{Process Management}
\version{2.0.0}

\begin{document}
    \begin{frame}[plain, noframenumbering]
        \titlepage
    \end{frame}
   

  \begin{slide}
    \slidetitle{The Operating System Creates Processes}

    The operating system has to:
    \begin{itemize}
      \item Maintain process control blocks, including state
      \item Create new processes
      \item Load a program, and re-initialize a process with context
    \end{itemize}

  \end{slide}
  
        \begin{slide}
    
        \slidetitle{On Unix, the Kernel Launches A Single User Process}

        After the kernel initializes, it creates a single process
        from a program
        \bigskip


        This process is called \texttt{init}, and it looks for it in \texttt{/sbin/init}

        \leftspace{}Responsible for executing every other process on the machine

        \leftspace{}Must always be active, if it exits the kernel thinks you're shutting down
        \bigskip

        For Linux, \texttt{init} will probably be \texttt{systemd} but there's other options
        \bigskip

        
    \end{slide}
    
    \begin{slide}

        \slidetitle{A Typical Process Tree}

        \vspace{-2mm}
        \begin{tikzpicture}[sibling distance=12em]
          \tikzstyle{every node}=[ellipse, draw, font=\scriptsize\ttfamily, align=center]
          \node {1: init}
            child { node {187: journald} }
            child { node {536: systemd --user}
              child { node [xshift=-6em] {597: gnome-shell}
                child { node { 946: firefox }
                  child { node [xshift=12em] {996: firefox} }
                  child { node [xshift=9em] {1089: firefox} }
                  child { node [xshift=6.35em] {1147: firefox} }
                }
              }
              child { node {1071: gnome-terminal-server}
                child { node [xshift=9em] {1086: zsh}
                  child { node {1240: htop}
                  }
                }
              }
            }
            child { node {200: udevd} };
        \end{tikzpicture}

    \end{slide}


    \begin{slide}

        \slidetitle{Processes Are Assigned a Process ID (pid) On Creation}

        The process ID is just a number, and is unique for every \textbf{active}
        process
        \bigskip

        On most Linux systems the maximum pid 32768, and 0 is reserved (invalid)
        \bigskip

        Eventually the kernel will recycle a pid, after the process dies, for a new process
        \bigskip

        Remember: each process has its own \textit{address space} (independent view of memory)

    \end{slide}

    \begin{slide}

        \slidetitle{Maintaining the Parent/Child Relationship}

        Previously, we made sure that our parent exited last (by using \texttt{sleep})
        \bigskip

        What happens if the parent process exits first, and no longer exists?

    \end{slide}

    \begin{slide}

        \slidetitle{The Parent Process is Responsible for Its Child}

        The operating system sets the exit status when a process terminates

        (the process terminates by calling \texttt{exit})

        \leftspace{}It can't remove its PCB yet
        \bigskip

        The minimum acknowledgment the parent has to do is read the child's exit status
        \bigskip

        There's two situations:
        \begin{enumerate}
          \item The child exits first (zombie process)
          \item The parent exits first (orphan process)
        \end{enumerate}

    \end{slide}

    \begin{slide}

        \slidetitle{A Zombie Process Waits for Its Parent to Read Its Exit Status}

        The process is terminated, but it hasn't been acknowledged
        \medskip

        A process may have an error in it, where it never reads the child's exit status
        \medskip

        The operating system can interrupt the parent process to acknowledge the child

        \leftspace{}It is just a suggestion and the parent is free to ignore it

        \leftspace{}\leftspace{}This is a basic form of IPC called a signal
        \bigskip

        The operating system has to keep a zombie process until it's acknowledged

        \leftspace{}If the parent ignores it, the zombie process needs to wait to
        be re-parented
        
    \end{slide}

    \begin{slide}

        \slidetitle{An Orphan Process Needs a New Parent}

        The child process lost its parent process

        \leftspace{}The child still needs a process to acknowledge its exit
        \medskip

        The operating system re-parents the child process to \texttt{init}

        \leftspace{}The \texttt{init} process is now responsible to acknowledge the
                     child
                     
    \end{slide}

  \begin{slide}

    \slidetitle{
      \texttt{\bfseries orphan-example.c} The Parent Exits Before the Child
    }

    \vspace{-5mm}
    \begin{minted}{c}
int main(int argc, char *argv[]) {
    pid_t pid = fork();
    if (pid == 0) {
        printf("Child parent pid: %d\n", getppid());
        sleep(2);
        printf("Child parent pid (after sleep): %d\n", getppid());
    }
    else {
        sleep(1);
    }
    return 0;
}

    \end{minted}

  \end{slide}

\begin{slide}

    \slidetitle{\texttt{\bfseries wait} delays execution}
    
    \texttt{wait} uses the following API:
    \begin{itemize}
      \item \texttt{status}: Address to store the wait status of the process
      \item Returns the process ID of child process

            \leftspace{}-1: on failure

            \leftspace{}>0: the child with a change
    \end{itemize}
    \bigskip

    The wait status contains a bunch of information, including the exit code

    \leftspace{}Use \texttt{man wait} to find all the macros to query wait status

    \leftspace{}\leftspace{}You can use \texttt{waitpid} to wait on a specific child process

\end{slide}

\begin{slide}

    \slidetitle{\texttt{\bfseries wait} Example}
    
\begin{minted}{c}
int main() {
    int pid = getpid();
    fork();
    if (getpid() == pid){
        wait(NULL);
        printf("Hello A\n");
    }
    else{
        printf("Hello B\n");
    }
    return 0;
}
\end{minted}
    
\end{slide}

    \begin{slide}

        \slidetitle{
          \texttt{\bfseries zombie-example.c} The Parent Monitors the Child To Check
          Its State
        }

        \vspace{-5mm}
        \begin{minted}{c}
pid_t pid = fork();

if (pid == 0) 
    sleep(2);

else {
    int ret;
    sleep(1);
    printf("Child process state: ");
    ret = print_state(pid);
    sleep(2);
    printf("Child process state: ");
    ret = print_state(pid);
}
        \end{minted}

    \end{slide}

\begin{slide}

    \slidetitle{You're Responsible for Managing Processes}

    The operating system maintains a strict parent/child relationship
    \medskip

    You should be able to identify (and prevent) the following:
    \begin{itemize}
        \item Zombie processes
        \item Orphan processes
    \end{itemize}
        
\end{slide}
   

\begin{slide}

    \slidetitle{Review Question}
    
    A process P1 runs on a Linux-like OS with a single core system.
    \medskip
    
    When P1 is executing, a disk interrupt occurs, causing P1 to go to kernel mode to service that interrupt. 
    \medskip
    
    The interrupt delivers all the disk blocks that unblock a process P2 (which blocked earlier on the disk read). 
    \medskip
    
    The interrupt service routine has completed execution fully, and the OS is just about to return back to the user mode of P1. 
    \medskip
    
    At this point in time, what are the states (ready/running/blocked) of processes P1 and P2?
    
\end{slide}

\begin{slide}

    \slidetitle{Review Question}
    
    A process P invokes the default wait system call. Which of the following scenarios block P 
    (i.e., P does not return immedaitely)?
    \medskip
    
    (a) P has no children.
    \medskip
    
    (b) P has one child that is still running.
    \medskip
    
    (c) P has one child that is a zombie.
    \medskip
    
    (d) P has two children, one is running and one is a zombie.
    
\end{slide}

\begin{slide}

    \slidetitle{Review Question}
    
    What sequence of system calls is executed when you run "sleep 100" and press Enter?
    \medskip
    
    (a) wait-exec-fork
    \medskip
    
    (b) exec-wait-fork
    \medskip
    
    (c) fork-exec-wait
    \medskip
    
    (d) wait-fork-exec

\end{slide}

\begin{slide}

    \slidetitle{Review Question}
    
    A process \texttt{P} runs the following code:
    \medskip
    
    \begin{minted}{c}
int ret1 = fork();
int ret2 = fork();
    \end{minted}
    
    How many direct children does P have?
    
\end{slide}

\begin{slide}

    \slidetitle{Review Question}

	What is the behavior of this program?
	
    \begin{minted}{c}
int main() {
    char *command = "/bin/ls"; 
    char *args[] = {"ls", NULL};   
    int rc = fork(); 
    if (rc == 0) {
        close(STDOUT_FILENO);
        int fd = open("foo.txt");
        execve(command, args, NULL);
        exit(0);
    } else
        wait(NULL);
    return 0;
}
    \end{minted}

\end{slide}

 \begin{slide}

    \slidetitle{Review Question}

	How many times is "Hello world!" printed?
	\medskip
	
    \inputminted{c}{fork_example_1_2.c}
    
\end{slide}

\begin{slide}

    \slidetitle{Review Question}

    Will it produce the same output each time it is run? If so, what is the output?
    
\medskip

\begin{minted}{c}
int main() {
    int i = 4;
    while (i != 0) {
        int pid = fork();
        if (pid == 0) {
            i--;
        } else {
            printf("%d\n", i);
            exit(0);
        }
    }
    return 0;
}
\end{minted}
\end{slide}

\begin{slide}

\slidetitle{Review Question}

    Will it produce the same output each time it is run? If so, what is the output?
\bigskip

    \begin{minted}{c}
int main() {
    int i = 4;
    while (i != 0) {
        int pid = fork();
        if (pid == 0) {
            i--;
        } else {
            waitpid(pid, NULL, 0);
            printf("%d\n", i);
            exit(0);
        }
    }
    return 0;
}
    \end{minted}

\end{slide}

\begin{slide}

    \slidetitle{Review Question}
    
    What is the output?
    \medskip
    
    \begin{minted}{c}
int x = 3; 
void main() {
    while (x > 0) {
        int ret = fork();
        char *args[] = {"a.out", NULL};
        if (ret == 0) {
            printf("x = %d\n", x);
            x--;
            execve("./a.out", args, NULL);
        } else {
            wait(NULL);
        }
    }
}
    \end{minted}
    
\end{slide}

\begin{slide}

    \slidetitle{Review Question}
    
    What is the output?
    \medskip
    
    \begin{minted}{c}
int x = 3; 
void main() 
{ 
    while(x > 0) { 
        int ret = fork(); 
        if(ret == 0) { 
            printf("x = %d\n", x); 
        } else 
            wait(NULL); 
        x--; 
    } 
}
    \end{minted}
    
\end{slide}

\end{document}
